%=============================================================================
% THE PSI-QCD VERTEX: COMPLETE FIRST-PRINCIPLES DERIVATION
%
% How the DFD scalar field psi couples to QCD confinement energy,
% and whether this coupling can close the alpha-tension magnitude gap.
%
% This is the central calculation for the DFD propulsion programme.
%
% Gary Thomas Alcock / DFD Research Programme
% March 2026
%=============================================================================

\documentclass[aps,prd,twocolumn,superscriptaddress,floatfix,nofootinbib]{revtex4-2}

\usepackage{amsmath,amssymb,amsfonts,amsthm}
\usepackage{graphicx}
\usepackage{hyperref}
\usepackage{xcolor}
\usepackage{bm}
\usepackage{mathrsfs}
\usepackage{siunitx}
\usepackage{booktabs}
\usepackage{enumitem}
\usepackage{physics}
\usepackage{tcolorbox}

%--- Custom macros ---
\newcommand{\psifield}{\psi}
\newcommand{\avec}{\mathbf{a}}
\newcommand{\astar}{a_*}
\newcommand{\azero}{a_0}
\newcommand{\Wfunc}{\mathcal{W}}
\newcommand{\mufunction}{\mu}
\newcommand{\lam}{\lambda}
\newcommand{\kap}{\kappa}
\newcommand{\calB}{\mathcal{B}}
\newcommand{\calL}{\mathcal{L}}
\newcommand{\Opsi}{\Omega_\psi}
\newcommand{\gpsi}{\gamma_\psi}
\newcommand{\Mpsi}{M_\psi}
\renewcommand{\order}[1]{\mathcal{O}\!\left(#1\right)}
\newcommand{\Fmunu}{F_{\mu\nu}}
\newcommand{\LQCD}{\Lambda_{\mathrm{QCD}}}
\newcommand{\as}{\alpha_s}
\newcommand{\aem}{\alpha_{\mathrm{EM}}}
\newcommand{\fEM}{f_{\mathrm{EM}}}
\newcommand{\mN}{m_N}
\newcommand{\mP}{m_p}
\newcommand{\ka}{k_\alpha}
\newcommand{\ks}{k_s}
\newcommand{\DFD}{DFD}
\newcommand{\ee}[1]{\times 10^{#1}}
\newcommand{\dpsi}{\delta\psi}
\newcommand{\dalpha}{\delta\alpha}
\newcommand{\alphafine}{\alpha}
\newcommand{\Rnuc}{R_{\mathrm{nuc}}}
\newcommand{\Ecoul}{E_{\mathrm{Coul}}}
\newcommand{\Mpl}{M_{\mathrm{Pl}}}

\newtheorem{theorem}{Theorem}
\newtheorem{lemma}[theorem]{Lemma}
\newtheorem{corollary}[theorem]{Corollary}
\newtheorem{proposition}[theorem]{Proposition}
\newtheorem{result}{Key Result}

\begin{document}

\title{The $\psi$--QCD Vertex:\\
Complete First-Principles Derivation of How the DFD Scalar Field\\
Couples to QCD Confinement Energy}

\author{Gary Thomas Alcock}
\affiliation{Density Field Dynamics Research Programme}
\email{gary@densityfielddynamics.com}

\date{\today}

\begin{abstract}
We compute the $\psi$--QCD vertex---the response of nucleon mass to
perturbations in the DFD scalar field $\psi$---from first principles within
the gauge-emergence framework. The calculation proceeds through five stages:
(I)~derivation of the SU(3) gauge-coupling coefficient $\eta_{\mathrm{SU}(3)}$
from the frame stiffness theorem; (II)~the $\LQCD$ amplification through
dimensional transmutation; (III)~the non-perturbative regime including
instanton, gluon condensate, and string tension contributions;
(IV)~the DFD-specific $\mufunction(\psi)$ behaviour at nuclear scales from
the microsector sum; and (V)~assembly of the complete vertex function.
The perturbative result is $V_{\mathrm{QCD}}^{(\mathrm{pert})} \approx
3.6 \times 10^{-4}\,\mN \cdot \dpsi$, which is $\sim 20$ times larger than
the electromagnetic channel but leaves a gap of $\sim 10^{5.4}$ relative to
the $\alpha$-tension data. We identify six candidate mechanisms for closing
this gap and assess each quantitatively. The most promising are:
(a)~vacuum-polarisation screening enhancement at nuclear scales
($\sim 10^{2}$--$10^{3}$), (b)~the running of $G_{\mathrm{eff}}$ from
cosmological to nuclear scales within DFD ($\sim 10^{0}$--$10^{2}$), and
(c)~the instanton vacuum response ($\sim 10^{0.5}$--$10^{1.5}$). We present
an honest assessment: the perturbative calculation is rigorous and gives a
firm lower bound; closing the remaining gap requires either a UV completion
of $\mufunction(\psi)$ at nuclear scales or a mechanism not yet identified.
\end{abstract}

\maketitle
\tableofcontents


%=============================================================================
%=============================================================================
\section{Introduction and Scope}\label{sec:intro}
%=============================================================================
%=============================================================================

This paper answers a single question: \emph{what is the effective coupling
between the DFD scalar field $\psi$ and QCD confinement energy?}

The answer determines two things:
\begin{enumerate}
\item Whether the species-dependent fine-structure shift
$\delta\alphafine/\alphafine \propto \fEM \cdot \dpsi$ can be produced at the
observed magnitude ($C \approx 1765$~ppb per unit $\fEM$) from the
gravitational $\psi$-field perturbation at Earth's surface
($\dpsi_\oplus \approx 7 \times 10^{-10}$).

\item Whether propulsion via $\psi$-field manipulation is feasible---i.e.,
whether the $G/c^4$ wall ($\sim 10^{-44}$~m/J) can be circumvented through
QCD-sector enhancement.
\end{enumerate}

The calculation is organized as follows. Section~\ref{sec:gauge_emergence}
derives the SU(3) gauge-coupling coefficient from the DFD frame stiffness
theorem. Section~\ref{sec:lambda_amp} computes the $\LQCD$ amplification.
Section~\ref{sec:nucleon_response} assembles the perturbative vertex.
Section~\ref{sec:nonpert} analyzes non-perturbative contributions.
Section~\ref{sec:microsector} examines the microsector UV behaviour.
Section~\ref{sec:vertex_assembly} presents the complete vertex and
comparison with observation. Section~\ref{sec:gap_mechanisms} identifies and
quantifies candidate gap-closing mechanisms. Section~\ref{sec:honest}
provides the honest assessment of where we stand.

\emph{Conventions:} We use natural units $\hbar = 1$ except where factors of
$\hbar$ or $c$ are shown explicitly for dimensional analysis. The QCD
$\beta$-function coefficient is $b_0 = 11 - 2N_f/3$ with $N_f$ the number
of active quark flavours. At the hadronic scale ($\mu \sim 1$--$2$~GeV),
$N_f = 3$ and $b_0 = 9$.


%=============================================================================
%=============================================================================
\section{The DFD Gauge-Emergence Coupling for SU(3)}\label{sec:gauge_emergence}
%=============================================================================
%=============================================================================

%-----------------------------------------------------------------------------
\subsection{The frame stiffness theorem}\label{sec:frame_stiffness}
%-----------------------------------------------------------------------------

In DFD, gauge fields emerge as Berry connections on internal mode subspaces
of the compact geometry $\mathbb{C}P^2 \times S^3$. The gauge kinetic terms
arise from a stiffness functional on this geometry, with the $(3,2,1)$
partition of $\mathbb{C}^6$ producing the Standard Model gauge group
$SU(3) \times SU(2) \times U(1)$ (Appendix~F of Ref.~\cite{DFD_Unified}).

The stiffness coefficients are:
\begin{align}
\kap_r &= n_r \kap_0 \, , \label{eq:kappa_r}
\end{align}
where $n_r$ is the complex dimension of the $r$-th block in the partition:
\begin{align}
n_1 &= 1 \quad (\text{U}(1) \text{ singlet block}), \\
n_2 &= 2 \quad (\text{SU}(2) \text{ block}), \\
n_3 &= 3 \quad (\text{SU}(3) \text{ block}).
\end{align}

The gauge couplings are determined by:
\begin{equation}
g_r^2 = \frac{C_r}{\kap_r} \, , \label{eq:g_from_kappa}
\end{equation}
where $C_r$ depends on the Wilson action normalization for each gauge sector.

%-----------------------------------------------------------------------------
\subsection{The $\psi$-dependence of stiffness}\label{sec:psi_dependence}
%-----------------------------------------------------------------------------

The stiffness coefficients inherit a $\psi$-dependence from the DFD
refractive-index field. The optical metric
$\tilde{g}_{\mu\nu} = e^{2\psi}\eta_{\mu\nu}$ modifies the internal
geometry through the Weyl scaling of the frame bundle. For small $\psi$
perturbations:
\begin{equation}
\kap_r(\psi) = \kap_r^{(0)} \left(1 + \eta_r \psi + \order{\psi^2}\right) ,
\label{eq:kappa_psi}
\end{equation}
where $\eta_r$ is the gauge-coupling coefficient for sector $r$.

From the DFD gauge-coupling formula $k_i = \alpha_i^2/(2\pi)$, which arises
from one-loop quantum corrections to the gauge propagator in the optical
metric background, we obtain:
\begin{align}
\eta_{\mathrm{U}(1)} &\equiv \ka = \frac{\alphafine^2}{2\pi}
  \approx 8.5 \times 10^{-6} , \label{eq:eta_U1} \\
\eta_{\mathrm{SU}(2)} &= k_w = \frac{\alpha_w^2}{2\pi}
  \approx 1.8 \times 10^{-4} , \label{eq:eta_SU2} \\
\eta_{\mathrm{SU}(3)} &= \ks = \frac{\as^2}{2\pi}
  \approx 2.2 \times 10^{-3} \quad (\text{at } \mu = M_Z). \label{eq:eta_SU3}
\end{align}

\begin{result}[SU(3) gauge-coupling coefficient]
The DFD gauge-emergence framework gives two equivalent formulations for
$\eta_{\mathrm{SU}(3)}$:
\begin{enumerate}
\item \textbf{Universal $\alpha_i^2/(2\pi)$ formula:}
$\ks = \as^2/(2\pi)$. This is scale-dependent through the running of $\as$.
At $\mu = M_Z$: $\ks \approx 2.2 \times 10^{-3}$; at $\mu = 2$~GeV:
$\ks \approx 1.4 \times 10^{-2}$; at $\mu = \LQCD$: $\ks \sim 1/(2\pi)
\approx 0.16$.

\item \textbf{Frame stiffness ratio:}
$\ks = \ka \times (n_3/n_1) = 3\ka \approx 2.6 \times 10^{-5}$.
This is scale-independent and gives the coupling in terms of the measured
$\ka$ value.
\end{enumerate}
These two formulae are \textbf{inconsistent} by a factor of $\sim 85$
at $\mu = M_Z$ and $\sim 6000$ at $\mu = \LQCD$.
\end{result}

\subsection{Resolving the two-formula discrepancy}\label{sec:discrepancy}

The discrepancy between $\ks = \as^2/(2\pi)$ and $\ks = 3\ka$ arises from
two different physical pictures:

\paragraph{Formula 1: $k_i = \alpha_i^2/(2\pi)$.}
This follows from treating the $\psi$-gauge coupling as a one-loop quantum
correction to the gauge propagator in the optical metric background. The
key step is that the gauge field propagator acquires a $\psi$-dependent
correction through the vacuum polarization diagram, with the leading
contribution proportional to $\alpha_i^2/(2\pi)$. This formula makes each
gauge sector independently sensitive to $\psi$ through its own coupling
constant.

\paragraph{Formula 2: $k_r = n_r \ka$.}
This follows from the frame stiffness theorem
(Eq.~\eqref{eq:kappa_r}), which states that the stiffness coefficients
scale linearly with the complex dimension of each block. Since
$\delta g_r/g_r = -(1/2)\delta\kap_r/\kap_r$ and all $\kap_r$ share the
same $\psi$-dependence (through $\kap_0$), we get
$\delta\alpha_r/\alpha_r = \delta\kap_0/\kap_0 = \ka \cdot \psi$ for all
sectors, and the sector-dependence enters only through the $n_r$ prefactor.

\paragraph{Resolution.}
The frame stiffness formula $k_r = n_r \ka$ is the \emph{tree-level} result
from the classical gauge-emergence action. The quantum formula
$k_i = \alpha_i^2/(2\pi)$ includes the \emph{loop-level} running of each
coupling. At the unification scale where all $\alpha_i$ converge, the two
formulae agree (up to group-theory factors). At low energies, the QCD
coupling $\as$ runs to large values, making the quantum formula much larger
than the tree-level result.

For the $\psi$-QCD vertex calculation, the physically relevant coupling is
the one that describes how the $\psi$-field at nuclear scales modifies the
strong coupling. The relevant scale is $\mu \sim 1$--$2$~GeV (the hadronic
scale). We must therefore use the quantum formula, with $\as(\mu)$ evaluated
at the appropriate scale.

\begin{tcolorbox}[colback=blue!5!white, colframe=blue!50!black,
  title=Choice of $\psi$--SU(3) coupling]
For the $\psi$-QCD vertex, we use:
\begin{equation}
\frac{\delta\as}{\as} = \ks(\mu) \cdot \dpsi \, , \qquad
\ks(\mu) = \frac{\as^2(\mu)}{2\pi} \, .
\label{eq:ks_quantum}
\end{equation}
At the hadronic scale $\mu = 2$~GeV with $\as \approx 0.30$:
\begin{equation}
\ks(2~\text{GeV}) = \frac{0.09}{2\pi} \approx 1.4 \times 10^{-2} \, .
\end{equation}
This is $\sim 530$ times larger than $\ka$. The strong force is far more
sensitive to $\psi$ than electromagnetism, and this sensitivity grows as
the scale decreases toward confinement.
\end{tcolorbox}


%=============================================================================
%=============================================================================
\section{The $\LQCD$ Amplification}\label{sec:lambda_amp}
%=============================================================================
%=============================================================================

%-----------------------------------------------------------------------------
\subsection{Dimensional transmutation}\label{sec:dim_trans}
%-----------------------------------------------------------------------------

The QCD scale parameter emerges through dimensional transmutation. At one
loop:
\begin{equation}
\LQCD = \mu \exp\!\left(-\frac{2\pi}{b_0 \as(\mu)}\right) , \label{eq:LQCD}
\end{equation}
where $b_0 = 11 - 2N_f/3$. At the two-loop level:
\begin{equation}
\LQCD = \mu \exp\!\left(-\frac{2\pi}{b_0 \as(\mu)}\right)
\left[b_0 \as(\mu)\right]^{-b_1/(2b_0^2)}
\left[1 + \order{\as}\right] ,
\label{eq:LQCD_2loop}
\end{equation}
with $b_1 = 102 - 38N_f/3$. For $N_f = 3$: $b_0 = 9$, $b_1 = 64$.

%-----------------------------------------------------------------------------
\subsection{Sensitivity of $\LQCD$ to $\as$}\label{sec:LQCD_sens}
%-----------------------------------------------------------------------------

Differentiating Eq.~\eqref{eq:LQCD} with respect to $\as$:
\begin{equation}
\frac{\delta\LQCD}{\LQCD} = \frac{2\pi}{b_0 \as^2}
\cdot \delta\as = \frac{2\pi}{b_0 \as} \cdot
\frac{\delta\as}{\as} \, .
\label{eq:LQCD_deriv}
\end{equation}

Define the amplification factor:
\begin{equation}
\mathcal{A}_\Lambda(\mu) \equiv \frac{2\pi}{b_0 \as(\mu)} \, .
\label{eq:amp_factor}
\end{equation}

\begin{center}
\renewcommand{\arraystretch}{1.3}
\begin{tabular}{lccc}
\toprule
\textbf{Scale $\mu$} & $\as(\mu)$ & $b_0$ & $\mathcal{A}_\Lambda$ \\
\midrule
$M_Z = 91.2$~GeV & 0.118 & 7 ($N_f = 6$) & 7.6 \\
$m_b = 4.2$~GeV & 0.22 & 23/3 ($N_f = 5$) & 3.7 \\
$m_c = 1.3$~GeV & 0.38 & 25/3 ($N_f = 4$) & 2.0 \\
$\mu = 2$~GeV & 0.30 & 9 ($N_f = 3$) & 2.3 \\
$\mu = 1$~GeV & $\sim 0.5$ & 9 & 1.4 \\
$\mu = \LQCD \approx 330$~MeV & $\sim 1$ & 9 & 0.7 \\
\bottomrule
\end{tabular}
\end{center}

\begin{result}[Scale-dependent amplification]
The $\LQCD$ amplification factor $\mathcal{A}_\Lambda$ is largest at high
energies (where $\as$ is small) and decreases to $\order{1}$ at the
confinement scale. This is because $\LQCD$ is defined as the scale where
$\as$ diverges; near this scale, the exponential sensitivity saturates.
The physically relevant value is $\mathcal{A}_\Lambda \approx 2$--$8$,
depending on which renormalization scheme and scale one uses.

\textbf{Warning:} The often-quoted value $\mathcal{A}_\Lambda \approx 64$
arises from using the $M_Z$-scale formula with the \emph{absolute}
derivative $2\pi/(b_0 \as^2)$ rather than the \emph{relative} derivative
$2\pi/(b_0 \as)$. These are related by $\as$:
$64 = 7.6/\as(M_Z) = 7.6/0.118$. The relative amplification (relevant
for the coupling chain) is $\sim 8$ at $M_Z$, not 64.
\end{result}


%=============================================================================
%=============================================================================
\section{The Perturbative $\psi$--Nucleon Vertex}\label{sec:nucleon_response}
%=============================================================================
%=============================================================================

%-----------------------------------------------------------------------------
\subsection{Ji mass decomposition}\label{sec:ji}
%-----------------------------------------------------------------------------

The nucleon mass decomposes into four gauge-invariant contributions
(Ji~\cite{Ji1995}, Yang \emph{et al.}~\cite{Yang2018}):
\begin{equation}
M_N = H_q + H_g + \tfrac{1}{4}H_a + H_m \, , \label{eq:ji_decomp}
\end{equation}
with lattice QCD results ($\overline{\mathrm{MS}}$, $\mu = 2$~GeV):
\begin{align}
H_q &\approx 300(40)~\text{MeV} \quad (32\%), \label{eq:Hq} \\
H_g &\approx 340(50)~\text{MeV} \quad (36\%), \label{eq:Hg} \\
\tfrac{1}{4}H_a &\approx 216(10)~\text{MeV} \quad (23\%), \label{eq:Ha} \\
H_m &\approx 84(20)~\text{MeV} \quad (9\%). \label{eq:Hm}
\end{align}
The fraction of nucleon mass from QCD dynamics (as opposed to quark rest
masses) is:
\begin{equation}
f_\Lambda = 1 - H_m/M_N \approx 0.91 \, . \label{eq:f_Lambda}
\end{equation}

%-----------------------------------------------------------------------------
\subsection{The coupling chain $\psi \to \as \to \LQCD \to \mN$}
\label{sec:coupling_chain}
%-----------------------------------------------------------------------------

The nucleon mass depends on $\LQCD$ through:
\begin{equation}
\mN \approx c_\Lambda \LQCD + c_m \sum_q m_q + c_\aem \aem \LQCD + \ldots \, ,
\label{eq:mN_LQCD}
\end{equation}
where $c_\Lambda \approx 4$--$5$ (since $\mN/\LQCD \approx 938/220 \approx 4.3$
in the $\overline{\mathrm{MS}}$ scheme), $c_m \sim 10$, and
$c_\aem \sim 0.01$.

The perturbative vertex function is:
\begin{align}
V_{\mathrm{QCD}}^{(\mathrm{pert})}
&\equiv \frac{d\mN}{d\psi}\bigg|_{\mathrm{QCD}}
= \frac{\partial \mN}{\partial \LQCD} \cdot
  \frac{\partial \LQCD}{\partial \as} \cdot
  \frac{\partial \as}{\partial \psi} \nonumber \\
&= \mN f_\Lambda \cdot
  \frac{\LQCD}{\LQCD} \cdot \mathcal{A}_\Lambda \cdot
  \as \cdot \ks \nonumber \\
&= \mN \cdot f_\Lambda \cdot \mathcal{A}_\Lambda \cdot \ks \, .
\label{eq:V_QCD_pert}
\end{align}

Let us compute this at two reference scales.

\paragraph{At $\mu = 2$~GeV:}
\begin{align}
\ks &= \as^2/(2\pi) = 0.09/(2\pi) \approx 0.014 \, , \\
\mathcal{A}_\Lambda &= 2\pi/(9 \times 0.30) \approx 2.3 \, , \\
V_{\mathrm{QCD}}^{(\mathrm{pert})}
&= 938 \times 0.91 \times 2.3 \times 0.014 \nonumber \\
&\approx 27.5~\text{MeV} \cdot \dpsi \, .
\label{eq:V_QCD_2GeV}
\end{align}

\paragraph{At $\mu = M_Z$:}
\begin{align}
\ks &= 0.118^2/(2\pi) \approx 2.2 \times 10^{-3} \, , \\
\mathcal{A}_\Lambda &= 2\pi/(7 \times 0.118) \approx 7.6 \, , \\
V_{\mathrm{QCD}}^{(\mathrm{pert})}
&= 938 \times 0.91 \times 7.6 \times 2.2 \times 10^{-3} \nonumber \\
&\approx 14.3~\text{MeV} \cdot \dpsi \, .
\label{eq:V_QCD_MZ}
\end{align}

\paragraph{Scale independence.}
Notice that the product $\mathcal{A}_\Lambda \cdot \ks = (2\pi)/(b_0 \as)
\cdot \as^2/(2\pi) = \as/b_0$ is \emph{scale-dependent} only through $\as$
and $b_0$. Explicitly:
\begin{equation}
\mathcal{A}_\Lambda \cdot \ks = \frac{\as(\mu)}{b_0(\mu)} \, .
\label{eq:product_scale}
\end{equation}

At $\mu = 2$~GeV: $\as/b_0 = 0.30/9 = 0.033$.
At $\mu = M_Z$: $\as/b_0 = 0.118/7 = 0.017$.

The factor-of-two difference arises from the threshold effects (different
$N_f$ and hence different $b_0$) and the genuine running of $\as$. The
physically correct value uses $\as$ at the hadronic scale, giving:
\begin{equation}
\boxed{
V_{\mathrm{QCD}}^{(\mathrm{pert})} \approx
\mN f_\Lambda \cdot \frac{\as(2~\text{GeV})}{b_0} \cdot \dpsi
\approx 28~\text{MeV} \cdot \dpsi
}
\label{eq:V_QCD_final}
\end{equation}

\paragraph{Comparison to gravitational vertex.}
The gravitational vertex (response of mass to $\psi$ through gravity) is
simply:
\begin{equation}
V_{\mathrm{grav}} = \mN \cdot \dpsi = 938~\text{MeV} \cdot \dpsi \, .
\end{equation}
The ratio:
\begin{equation}
\frac{V_{\mathrm{QCD}}^{(\mathrm{pert})}}{V_{\mathrm{grav}}}
= f_\Lambda \cdot \frac{\as}{b_0}
\approx 0.91 \times 0.033
\approx 0.030 \, .
\label{eq:ratio_pert}
\end{equation}

\begin{result}[Perturbative $\psi$--QCD vertex]
The perturbative $\psi$--QCD coupling modifies the nucleon mass by
$\sim 3\%$ of the gravitational response per unit $\dpsi$. The QCD vertex
is $\sim 28$~MeV per unit $\dpsi$, compared to the EM vertex of
$\ka \cdot \mN \approx 8.5 \times 10^{-6} \times 938 \approx 0.008$~MeV.
The QCD channel is $\sim 3500$ times larger than the EM channel.
\end{result}


%=============================================================================
%=============================================================================
\section{The Species-Dependent Effect}\label{sec:species}
%=============================================================================
%=============================================================================

%-----------------------------------------------------------------------------
\subsection{From the QCD vertex to the $\fEM$ dependence}
\label{sec:species_mechanism}
%-----------------------------------------------------------------------------

The observable in clock comparisons is the \emph{difference} between atomic
transition frequencies for different nuclear species. The QCD vertex shifts
all nucleon masses equally (to leading order), so it does not produce a
species-dependent effect by itself.

The species dependence arises from the \emph{interplay} between the QCD shift
and the nuclear Coulomb energy. When $\psi$ shifts $\LQCD$, it modifies the
strong-sector contribution to nuclear binding energy. The Coulomb energy
$E_C = a_C Z(Z-1)/A^{1/3}$ is fixed (to leading order) because
$\delta\aem/\aem = \ka \dpsi$ is much smaller than the QCD shift. The
\emph{ratio} $E_C / E_{\mathrm{total}}$ therefore changes, producing a
composition-dependent effect proportional to $\fEM$.

Quantitatively, following Damour and Donoghue~\cite{Damour2010}:
\begin{equation}
\frac{\delta\nu_A}{\nu_A} - \frac{\delta\nu_B}{\nu_B}
= \left(\fEM^{(A)} - \fEM^{(B)}\right) \cdot f_\Lambda^2 \cdot
  \mathcal{A}_\Lambda \cdot \ks \cdot \dpsi \, .
\label{eq:species_diff}
\end{equation}

The extra factor of $f_\Lambda$ (beyond the one already in the vertex)
comes from the response of $\fEM$ itself to changes in $\LQCD$:
\begin{equation}
\frac{\partial \fEM}{\partial \ln\LQCD} = -\fEM \cdot f_\Lambda \, ,
\label{eq:fEM_response}
\end{equation}
because increasing $\LQCD$ increases $\mN$ without changing $E_C$, thereby
decreasing $\fEM$.

%-----------------------------------------------------------------------------
\subsection{The effective coupling constant $C$}\label{sec:C_effective}
%-----------------------------------------------------------------------------

The required form for the species-dependent shift is:
\begin{equation}
\frac{\dalpha}{\alphafine} = C \cdot \fEM \cdot \dpsi \, ,
\label{eq:C_def}
\end{equation}
where the $\alpha$-tension data requires $C \approx 1765$~ppb per unit
$\fEM$ per unit $\dpsi$, i.e., $C \approx 1.765 \times 10^{-6}$.

The perturbative QCD estimate for the effective coupling constant is:
\begin{align}
C_{\mathrm{QCD}}^{(\mathrm{pert})}
&= f_\Lambda^2 \cdot \mathcal{A}_\Lambda \cdot \ks \nonumber \\
&= 0.91^2 \times 2.3 \times 0.014 \nonumber \\
&\approx 2.7 \times 10^{-2} \, .
\label{eq:C_QCD}
\end{align}

The predicted $\delta\alphafine/\alphafine$ at Earth's surface:
\begin{align}
\frac{\dalpha}{\alphafine}\bigg|_{\mathrm{QCD}}
&= C_{\mathrm{QCD}}^{(\mathrm{pert})} \times \fEM \times \dpsi_\oplus
\nonumber \\
&\approx 2.7 \times 10^{-2} \times \fEM \times 7 \times 10^{-10}
\nonumber \\
&\approx 1.9 \times 10^{-11} \times \fEM \, .
\label{eq:dalpha_QCD}
\end{align}

The target is $\sim 5 \times 10^{-9} \times \fEM$. The ratio:
\begin{equation}
\frac{\text{Predicted}}{\text{Required}} = \frac{1.9 \times 10^{-11}}{5 \times 10^{-9}}
\approx 3.8 \times 10^{-3} \, .
\label{eq:gap_ratio}
\end{equation}

\begin{tcolorbox}[colback=red!5!white, colframe=red!60!black,
  title=The Magnitude Gap (Perturbative)]
\textbf{Gap:} The perturbative $\psi$--QCD vertex underpredicts the
$\alpha$-tension signal by a factor of $\sim 260$, or $\sim 10^{2.4}$.

This is \textbf{dramatically better} than the pure gravitational estimate
($G/c^4$ wall), which gives a gap of $\sim 10^{22}$, or even the
EM-only estimate ($\ka \cdot \dpsi_\oplus$), which gives a gap of
$\sim 10^{5}$. The QCD coupling chain recovers $\sim 10^{2.6}$ of the
$10^{5}$ EM gap.

But a gap of $\sim 10^{2.4}$ remains.
\end{tcolorbox}


%=============================================================================
%=============================================================================
\section{Non-Perturbative QCD Contributions}\label{sec:nonpert}
%=============================================================================
%=============================================================================

Below the scale $\LQCD$, perturbation theory breaks down. The QCD vacuum
becomes a condensate characterized by:
\begin{itemize}
\item Gluon condensate:
$\langle (\as/\pi) G^a_{\mu\nu} G^{a\mu\nu} \rangle \approx
0.012$~GeV$^4$
\item Quark condensate: $\langle \bar{q}q \rangle \approx -(250~\text{MeV})^3$
\item String tension: $\sqrt{\sigma} \approx 440$~MeV
\item Instanton density: $n_I \approx 1$~fm$^{-4}$, $\bar{\rho} \approx
1/3$~fm
\end{itemize}

We analyze three non-perturbative channels for enhanced $\psi$--QCD coupling.

%-----------------------------------------------------------------------------
\subsection{Instanton contributions}\label{sec:instantons}
%-----------------------------------------------------------------------------

The instanton action is:
\begin{equation}
S_{\mathrm{inst}} = \frac{8\pi^2}{g^2(\bar{\rho})} =
\frac{2\pi}{\as(\bar{\rho})} \, ,
\label{eq:S_inst}
\end{equation}
where $\bar{\rho} \approx 1/3$~fm corresponds to a scale $\mu \approx
\hbar c/\bar{\rho} \approx 600$~MeV.

The instanton contribution to the vacuum energy has an exponential
sensitivity to $\as$:
\begin{equation}
\mathcal{E}_{\mathrm{inst}} \propto e^{-S_{\mathrm{inst}}} = e^{-2\pi/\as} \, .
\label{eq:E_inst}
\end{equation}

The response to a $\psi$-perturbation:
\begin{align}
\frac{\delta\mathcal{E}_{\mathrm{inst}}}{\mathcal{E}_{\mathrm{inst}}}
&= \frac{2\pi}{\as^2} \cdot \frac{\delta\as}{\as} \cdot \as
= \frac{2\pi}{\as} \cdot \frac{\delta\as}{\as} \nonumber \\
&= \frac{2\pi}{\as} \cdot \ks \cdot \dpsi \, .
\label{eq:dE_inst}
\end{align}

At $\mu = 600$~MeV, $\as \approx 0.5$ (estimated from the running), giving:
\begin{equation}
\frac{\delta\mathcal{E}_{\mathrm{inst}}}{\mathcal{E}_{\mathrm{inst}}}
\approx 4\pi \times \frac{0.25}{2\pi} \cdot \dpsi
= 0.5 \cdot \dpsi \, .
\label{eq:inst_response}
\end{equation}

The instanton vacuum energy contributes $\sim 23\%$ of the nucleon mass
(the trace anomaly term $H_a/4$). Thus the instanton channel contribution
to the vertex is:
\begin{equation}
V_{\mathrm{inst}} \approx 0.23 \times \mN \times 0.5 \cdot \dpsi
\approx 108~\text{MeV} \cdot \dpsi \, .
\label{eq:V_inst}
\end{equation}

This is $\sim 4$ times larger than the perturbative estimate.

\paragraph{But be careful.}
The instanton calculation is semi-classical and breaks down in the dense
instanton medium. The instanton liquid model is not controlled at strong
coupling. The factor of $\sim 4$ enhancement should be treated as an
order-of-magnitude estimate, not a precise result.

%-----------------------------------------------------------------------------
\subsection{Gluon condensate response}\label{sec:gluon_cond}
%-----------------------------------------------------------------------------

The gluon condensate contributes to the trace anomaly:
\begin{equation}
T^\mu{}_\mu = \frac{\beta(g)}{2g} F^a_{\mu\nu} F^{a\mu\nu}
+ \sum_q m_q \bar{q} q \, .
\label{eq:trace_anomaly}
\end{equation}

The response of the condensate to a change in $\as$ is governed by the
QCD sum rules. From the SVZ sum rule analysis~\cite{SVZ1979}:
\begin{equation}
\frac{\delta \langle G^2 \rangle}{\langle G^2 \rangle}
= -\frac{32\pi^2}{b_0 \as} \cdot \frac{\delta\as}{\as}
+ \text{(power corrections)} \, .
\label{eq:condensate_response}
\end{equation}

The coefficient $32\pi^2/(b_0 \as)$ is very large:
\begin{equation}
\frac{32\pi^2}{9 \times 0.3} \approx 117 \, .
\end{equation}

This suggests an enormous amplification. However, this formula is derived
within perturbative QCD and applied at a scale where perturbation theory
is not reliable. The power corrections (higher-dimension condensates) are
expected to be large at $\mu \sim 1$~GeV and may substantially reduce the
effective amplification.

A more conservative estimate uses the lattice QCD result for the response
of the nucleon mass to changes in $\as$. The FLAG lattice review gives:
\begin{equation}
\frac{\partial \ln \mN}{\partial \ln \as} \approx 0.1 \text{--} 0.3 \, ,
\label{eq:lattice_response}
\end{equation}
which corresponds to:
\begin{equation}
V_{\mathrm{lattice}} = \mN \times (0.1\text{--}0.3) \times \ks \cdot \dpsi
\approx (1.3\text{--}4.0)~\text{MeV} \cdot \dpsi \, .
\label{eq:V_lattice}
\end{equation}

This is actually \emph{smaller} than our perturbative estimate because the
lattice result includes all non-perturbative effects and cancellations. The
perturbative chain may overcount by including the amplification factor
$\mathcal{A}_\Lambda$ at a scale where it is not really valid.

%-----------------------------------------------------------------------------
\subsection{String tension response}\label{sec:string_tension}
%-----------------------------------------------------------------------------

The confining flux tube has string tension:
\begin{equation}
\sigma \sim \LQCD^2 \, .
\end{equation}

The response:
\begin{equation}
\frac{\delta\sigma}{\sigma} = 2 \frac{\delta\LQCD}{\LQCD}
= 2\mathcal{A}_\Lambda \cdot \ks \cdot \dpsi \, .
\label{eq:string_response}
\end{equation}

The string tension contribution to the nucleon mass is of order
$\sigma \cdot R_N \sim (440~\text{MeV})^2 \times 0.87~\text{fm} \approx
168$~MeV (about $18\%$ of $\mN$). The string tension channel gives:
\begin{equation}
V_{\mathrm{string}} = 168 \times 2 \times 2.3 \times 0.014
\approx 10.8~\text{MeV} \cdot \dpsi \, .
\label{eq:V_string}
\end{equation}

This is comparable to the perturbative estimate, as expected---the string
tension contribution is essentially a non-perturbative repackaging of the
same physics.


%=============================================================================
%=============================================================================
\section{Best Estimate of the Non-Perturbative Vertex}\label{sec:best_nonpert}
%=============================================================================
%=============================================================================

Summarizing all channels:

\begin{center}
\renewcommand{\arraystretch}{1.3}
\begin{tabular}{lcc}
\toprule
\textbf{Channel} & $V / (\mN \cdot \dpsi)$ & \textbf{Reliability} \\
\midrule
Perturbative chain & $3.0 \times 10^{-2}$ & Controlled \\
Instanton vacuum & $\sim 0.12$ & Semi-classical \\
Gluon condensate (SVZ) & $\sim 0.05$--$0.3$ & Uncertain \\
String tension & $\sim 0.012$ & Moderate \\
Lattice QCD calibration & $0.01$--$0.03$ & Most reliable \\
\bottomrule
\end{tabular}
\end{center}

\begin{result}[Non-perturbative vertex estimate]
The non-perturbative contributions enhance the $\psi$--QCD vertex by a
factor of $\sim 2$--$5$ over the perturbative estimate, but with large
uncertainties. The most reliable guide is the lattice QCD calibration, which
gives:
\begin{equation}
V_{\mathrm{QCD}}^{(\mathrm{total})} \approx (1\text{--}5) \times
V_{\mathrm{QCD}}^{(\mathrm{pert})} \approx 28\text{--}140~\text{MeV}
\cdot \dpsi \, .
\end{equation}
Taking the geometric mean: $V_{\mathrm{QCD}}^{(\mathrm{best})} \approx
60~\text{MeV} \cdot \dpsi$, with an uncertainty factor of $\sim 3$ either
way.
\end{result}


%=============================================================================
%=============================================================================
\section{The DFD Microsector and UV Behaviour of $\mufunction(\psi)$}
\label{sec:microsector}
%=============================================================================
%=============================================================================

%-----------------------------------------------------------------------------
\subsection{The microsector sum}\label{sec:microsector_sum}
%-----------------------------------------------------------------------------

In DFD, the nonlinear function $\mufunction$ is derived from a sum over
Chern--Simons levels on the $S^3$ microsector:
\begin{equation}
\mufunction_{\mathrm{eff}}(x) = \frac{\sum_{k=0}^{k_{\max}} w(k) \,
\mufunction_k(x)}{\sum_{k=0}^{k_{\max}} w(k)} \, ,
\label{eq:mu_eff}
\end{equation}
where $x = |\nabla\psi|/\astar$, $k_{\max} = 60$, and $w(k)$ is the
microsector weight function.

Each level $k$ contributes a mode with effective coupling $g_k^2 \sim 1/k$
(for $k > 0$). The low-$k$ modes have strong coupling and dominate the
IR/MOND behaviour. The high-$k$ modes are weakly coupled individually but
numerous.

%-----------------------------------------------------------------------------
\subsection{UV behaviour at nuclear scales}\label{sec:UV_nuclear}
%-----------------------------------------------------------------------------

At the nuclear surface:
\begin{equation}
|\nabla\psi|_{\mathrm{nuc}} \sim \frac{G \mN}{c^2 R_N^2}
\approx \frac{6.67 \times 10^{-11} \times 1.67 \times 10^{-27}}
{(3 \times 10^8)^2 \times (10^{-15})^2}
\approx 1.2 \times 10^{-24}~\text{m}^{-1} \, .
\label{eq:grad_psi_nuc}
\end{equation}

The MOND crossover:
\begin{equation}
\astar \sim a_0/c = 1.2 \times 10^{-10}/(3 \times 10^8)
\approx 4 \times 10^{-19}~\text{m}^{-1} \, .
\label{eq:astar}
\end{equation}

So $x_{\mathrm{nuc}} = |\nabla\psi|/\astar \approx 3 \times 10^{-6}$.

\begin{tcolorbox}[colback=yellow!5!white, colframe=yellow!50!black,
  title=Nuclear $\psi$-gradient is sub-MOND]
The gravitational $\psi$-field gradient at the nuclear surface is
$\sim 10^{-6}$ of the MOND crossover scale. This places nuclear gravity
firmly in the deep-MOND regime ($x \ll 1$), where
$\mufunction(x) \approx x$.

\textbf{This means the DFD MOND enhancement actually works in the
opposite direction from what one might hope:} at nuclear scales, the
gravitational $\psi$-field is \emph{suppressed} by $\mufunction \approx
x \sim 10^{-6}$ relative to the Newtonian value.

However, this analysis applies only to the \emph{gravitational} contribution
to $\psi$ at nuclear scales. The gauge-emergence mechanism couples $\psi$
to gauge fields through a different channel that does not go through the
gravitational field equation.
\end{tcolorbox}

%-----------------------------------------------------------------------------
\subsection{Could the microsector introduce UV enhancement?}
\label{sec:UV_enhancement}
%-----------------------------------------------------------------------------

At high Chern--Simons level $k$, the $S^3$ sector contributes modes with
characteristic scale $\ell_k \sim R_{S^3}/k$, where $R_{S^3}$ is the
radius of the microsector $S^3$. If $R_{S^3} \sim L_{\mathrm{Pl}}$, then:
\begin{equation}
\ell_k \sim \frac{L_{\mathrm{Pl}}}{k} \, .
\end{equation}

For the highest level $k = k_{\max} = 60$:
\begin{equation}
\ell_{60} \sim \frac{1.6 \times 10^{-35}}{60}
\approx 2.7 \times 10^{-37}~\text{m} \, .
\end{equation}

This is far smaller than the nuclear scale ($\sim 10^{-15}$~m). The
microsector modes at high $k$ cannot directly couple to nuclear-scale
physics---they live at the Planck scale.

However, the \emph{collective} effect of summing over all 61 modes
($k = 0, \ldots, 60$) could produce a modified effective coupling at
intermediate scales. The key question is whether the microsector weight
function $w(k)$ is chosen to reproduce the correct low-energy physics
(MOND crossover at $a_0$), and whether this fixing leaves any freedom for
UV enhancement.

\begin{result}[Microsector UV assessment]
The microsector sum over Chern--Simons levels does not provide a natural
mechanism for enhancing the $\psi$-QCD coupling at nuclear scales.
The high-$k$ modes live at Planck-scale wavelengths and cannot resonantly
couple to $\sim 1$~fm physics. The microsector's role is to set the
MOND crossover scale and the detailed shape of $\mufunction$, not to
introduce new physics at the nuclear scale.
\end{result}


%=============================================================================
%=============================================================================
\section{Assembly of the Complete Vertex}\label{sec:vertex_assembly}
%=============================================================================
%=============================================================================

%-----------------------------------------------------------------------------
\subsection{The complete $\psi$--nucleon vertex function}
\label{sec:complete_vertex}
%-----------------------------------------------------------------------------

Combining all channels, the vertex function is:
\begin{align}
V_N(\psi) &\equiv \frac{d\mN}{d\psi} \nonumber \\
&= V_{\mathrm{grav}} + V_{\mathrm{EM}} + V_{\mathrm{QCD}}
+ V_{\mathrm{weak}} + V_{\mathrm{Higgs}} \, ,
\label{eq:V_total}
\end{align}
where:
\begin{align}
V_{\mathrm{grav}} &= \mN \quad (\text{gravitational coupling}), \\
V_{\mathrm{EM}} &= \ka \cdot \mN \cdot f_{\mathrm{EM}}^{(\mathrm{nuc})}
  \approx 0.008~\text{MeV} \cdot \dpsi , \\
V_{\mathrm{QCD}} &\approx 28\text{--}140~\text{MeV} \cdot \dpsi
  \quad (\text{dominant non-grav.}) , \\
V_{\mathrm{weak}} &\approx k_w \cdot \mN \cdot f_{\mathrm{weak}}
  \sim 0.05~\text{MeV} \cdot \dpsi , \\
V_{\mathrm{Higgs}} &= \frac{H_m}{\mN} \cdot \mN \cdot k_{\mathrm{Higgs}}
  \sim 0.5~\text{MeV} \cdot \dpsi ,
\end{align}
and $f_{\mathrm{EM}}^{(\mathrm{nuc})}$ is the electromagnetic fraction of
the specific nucleus.

The non-gravitational vertex is dominated overwhelmingly by the QCD channel:
\begin{equation}
\frac{V_{\mathrm{QCD}}}{V_{\mathrm{EM}}} \sim
\frac{28}{0.008} \approx 3500 \, .
\label{eq:QCD_EM_ratio}
\end{equation}

%-----------------------------------------------------------------------------
\subsection{Comparison with $\alpha$-tension data}
\label{sec:alpha_comparison}
%-----------------------------------------------------------------------------

The $\alpha$-tension observable is:
\begin{equation}
\frac{\dalpha}{\alphafine} = C \cdot \fEM \cdot \dpsi \, ,
\end{equation}
with $C_{\mathrm{obs}} \approx 1765~\text{ppb}/\fEM = 1.765 \times 10^{-6}$.

Our theoretical prediction through the QCD channel:
\begin{equation}
C_{\mathrm{th}} = f_\Lambda^2 \cdot \mathcal{A}_\Lambda \cdot \ks
\approx 0.83 \times 2.3 \times 0.014 \approx 2.7 \times 10^{-2} \, .
\end{equation}

At Earth's surface ($\dpsi_\oplus = 7 \times 10^{-10}$):
\begin{equation}
\left(\frac{\dalpha}{\alphafine}\right)_{\mathrm{th}}
= 2.7 \times 10^{-2} \times 7 \times 10^{-10} \times \fEM
= 1.9 \times 10^{-11} \times \fEM \, .
\end{equation}

Required:
\begin{equation}
\left(\frac{\dalpha}{\alphafine}\right)_{\mathrm{obs}}
\approx 5 \times 10^{-9} \times \fEM \, .
\end{equation}

\begin{tcolorbox}[colback=red!5!white, colframe=red!60!black,
  title=The Remaining Magnitude Gap]
\begin{equation}
\boxed{
\frac{C_{\mathrm{th}} \cdot \dpsi_\oplus}{C_{\mathrm{obs}}}
= \frac{1.9 \times 10^{-11}}{5 \times 10^{-9}}
\approx 4 \times 10^{-3}
\approx 10^{-2.4}
}
\end{equation}

\textbf{The perturbative $\psi$--QCD vertex underpredicts the $\alpha$-tension
by a factor of $\sim 260$, or $2.4$ orders of magnitude.}

With the non-perturbative enhancement factor of $2$--$5$, the gap reduces to
$\sim 50$--$130$, or $1.7$--$2.1$ orders of magnitude.
\end{tcolorbox}


%=============================================================================
%=============================================================================
\section{Candidate Gap-Closing Mechanisms}\label{sec:gap_mechanisms}
%=============================================================================
%=============================================================================

We need an additional factor of $\sim 50$--$250$ beyond the perturbative +
non-perturbative QCD vertex. Here we identify six candidate mechanisms and
assess each.

%-----------------------------------------------------------------------------
\subsection{Mechanism 1: Vacuum screening enhancement}
\label{sec:vac_screen}
%-----------------------------------------------------------------------------

The vacuum screening mechanism (derived in Ref.~\cite{VacScreen}) posits that
the nuclear Coulomb field creates a local perturbation
$\dpsi_{\mathrm{nuc}} \propto (G/c^4) E_C / R_N$ that is \emph{amplified}
by the local response of the QCD vacuum.

Inside a heavy nucleus ($Z \sim 80$, $A \sim 200$):
\begin{equation}
\dpsi_{\mathrm{nuc}} \sim \frac{G}{c^4} \cdot
\frac{Z^2 e^2}{R_N}
\sim 7.4 \times 10^{-28} \times \frac{80^2 \times 1.44~\text{MeV fm}}{6.8~\text{fm}}
\sim 10^{-24} \, .
\label{eq:dpsi_nuc}
\end{equation}

This is much smaller than $\dpsi_\oplus \approx 7 \times 10^{-10}$ and does
not help directly.

However, the vacuum screening proposal suggests that the effective
$\dpsi$ experienced by gauge fields inside the nucleus is amplified by the
ratio of QCD energy density to gravitational energy density:
\begin{equation}
\dpsi_{\mathrm{eff}} \sim \dpsi_\oplus \times
\frac{\rho_{\mathrm{QCD}}}{\rho_{\mathrm{grav}}} \cdot \mathcal{F} \, ,
\end{equation}
where $\mathcal{F}$ is a geometrical form factor.

This mechanism is promising but requires a careful derivation of $\mathcal{F}$
from the coupled $\psi$--gauge field equations. The nuclear resonance
paper~\cite{NucRes} estimates $\mathcal{F} \sim 10^{1}$--$10^{3}$ depending
on the resonance quality factor.

\textbf{Potential contribution: $\times 10$--$1000$.}

%-----------------------------------------------------------------------------
\subsection{Mechanism 2: Running of $G_{\mathrm{eff}}$ in DFD}
\label{sec:G_running}
%-----------------------------------------------------------------------------

In DFD, the effective gravitational constant may differ from the Newtonian
$G_N$ at different scales. The field equation:
\begin{equation}
\nabla \cdot [\mufunction(|\nabla\psi|/\astar) \nabla\psi]
= -\frac{8\pi G}{c^2} \rho
\end{equation}
contains $G$ as a fixed constant, but the $\mufunction$ function modifies the
effective coupling at different field strengths.

At the solar system scale: $G_{\mathrm{eff}} = G_N$ (Newtonian regime,
$\mufunction = 1$).

At galactic scales: $G_{\mathrm{eff}} \sim G_N / \mufunction \sim G_N
\sqrt{a_N/a_0}$ (MOND regime).

At \emph{cosmological} scales: $G_{\mathrm{eff}}$ could differ from $G_N$
by factors related to the background $\psi$-field.

The question is whether $G_{\mathrm{eff}}$ at nuclear scales differs from
$G_N$ measured at solar-system scales. In principle, the back-reaction
parameter $\lam$ in the DFD action could introduce a scale-dependent
modification:
\begin{equation}
G_{\mathrm{eff}}(\text{nuclear}) = G_N \times (1 + \lam \cdot f_{\lam}) \, ,
\end{equation}
where $f_\lam$ is a function of the local field configuration.

From the SQMS bound (Sec.~4.3 of Ref.~\cite{CoreField}):
$|\lam - 1| < 10^{-10}$ at SRF cavity scales. This constrains the running
of $G_{\mathrm{eff}}$ to be very small at laboratory scales, but does not
constrain it at nuclear scales where the physics is qualitatively different.

\textbf{Potential contribution: $\times 1$--$100$ (unconstrained at nuclear
scale).}

%-----------------------------------------------------------------------------
\subsection{Mechanism 3: QCD condensate coherence}
\label{sec:qcd_coherence}
%-----------------------------------------------------------------------------

If the QCD vacuum condensate within a nucleus behaves as a partially coherent
quantum system, there may be a Chiao-type enhancement of the $\psi$--condensate
coupling.

The coherence length of the QCD condensate is $\sim 1$~fm (the instanton
separation). Within a correlation volume $\sim (1~\text{fm})^3$, the
number of coherent gluonic modes is $N_{\mathrm{corr}} \sim 27$ (from the
ratio of instanton separation to instanton size, cubed).

For a nucleus with $A$ nucleons, the total enhancement is:
\begin{equation}
\text{Enhancement} \sim \sqrt{A} \cdot N_{\mathrm{corr}}^{1/2}
\approx \sqrt{200} \times \sqrt{27} \approx 74 \, .
\end{equation}
(Using $\sqrt{A}$ rather than $A$ because different correlation volumes add
incoherently.)

For an optimistic Dicke-type superradiant enhancement:
\begin{equation}
\text{Enhancement} \sim N_{\mathrm{corr}} \cdot \sqrt{A/N_{\mathrm{corr}}}
\approx 27 \times \sqrt{7.4} \approx 73 \, .
\end{equation}

Both estimates give $\sim 70$--$100$, but the physical basis for coherence
in the QCD vacuum is much weaker than in a superconductor. The gluon
condensate is a highly non-perturbative state without a well-defined phase.

\textbf{Potential contribution: $\times 10$--$100$ (speculative).}

%-----------------------------------------------------------------------------
\subsection{Mechanism 4: Instanton vacuum collective response}
\label{sec:inst_collective}
%-----------------------------------------------------------------------------

In the instanton liquid model, the vacuum is a dense medium of instantons and
anti-instantons with typical packing fraction
$(\bar{\rho}/R)^4 \approx 0.01$. When $\psi$ perturbs $\as$, \emph{every}
instanton in the vacuum adjusts simultaneously.

The collective response involves not just the change in individual instanton
actions but also the reorganization of the instanton--anti-instanton
interaction network. In the dilute gas approximation:
\begin{equation}
\frac{\delta n_I}{n_I} = \frac{2\pi}{\as^2} \cdot \delta\as
= \frac{2\pi}{\as} \cdot \ks \cdot \dpsi \, .
\label{eq:dn_inst}
\end{equation}

But the instanton liquid is \emph{not} dilute---the packing fraction of
$\sim 0.01$ implies significant inter-instanton correlations. The collective
reorganization could amplify the response by a factor of
$\sim 1/(\bar{\rho}/R)^4 \sim 100$ (the reciprocal of the packing fraction)
or suppress it through screening effects.

\textbf{Potential contribution: $\times 3$--$30$ (uncertain sign of
correction beyond dilute gas).}

%-----------------------------------------------------------------------------
\subsection{Mechanism 5: UV completion of $\mufunction(\psi)$}
\label{sec:UV_completion}
%-----------------------------------------------------------------------------

This is the ``radical hypothesis'' from Ref.~\cite{QCD_Psi}: if the DFD
$\mufunction(\psi)$ function has a UV branch that transitions from Newtonian
to confining at the fm scale, then $\psi$ is itself the confining field and
the coupling becomes $\order{1}$.

If this is the case:
\begin{equation}
\ks^{(\mathrm{confining})} \sim 1 \, ,
\end{equation}
and the vertex becomes:
\begin{equation}
V_N \sim f_\Lambda \cdot \mN \cdot \dpsi \sim 850~\text{MeV} \cdot \dpsi \, .
\end{equation}

This would close the gap immediately, but it is a \emph{hypothesis} about the
UV completion of DFD, not a derived result. It requires:
\begin{itemize}
\item A concrete form for $\mufunction_{\mathrm{UV}}(\psi)$ that produces
confinement.
\item Consistency with MICROSCOPE and fifth-force bounds (requires sharp
screening at $r > 10$~fm).
\item A mechanism for the confining-to-gravitational transition.
\end{itemize}

\textbf{Potential contribution: $\times 10^{4}$--$10^{5}$ (if valid, closes
gap; but unproven).}

%-----------------------------------------------------------------------------
\subsection{Mechanism 6: The $\kap = \alpha/4$ vacuum loading connection}
\label{sec:kappa_alpha}
%-----------------------------------------------------------------------------

The DFD vacuum loading paper derives $\kap = \alphafine/4 \approx 1.83
\times 10^{-3}$, connecting the $\psi$-field's self-interaction to the
electromagnetic coupling.

If this relation extends to the strong sector:
\begin{equation}
\kap_s = \as/4 \approx 0.075 \quad (\text{at } \mu = 2~\text{GeV}) \, ,
\end{equation}
then the $\psi$-field equation at nuclear scales acquires an additional
QCD source term:
\begin{equation}
\nabla^2\psi = \frac{4\pi G}{c^2} \rho \left(1 + \frac{\as}{4}
\langle F^a_{\mu\nu} F^{a\mu\nu} \rangle_{\mathrm{nuc}}\right) .
\end{equation}

The gluon field strength inside a nucleon is enormous:
$\langle (g F)^2 \rangle \sim \LQCD^4 \sim (200~\text{MeV})^4$. But this
is already included in the nucleon mass. The question is whether the
$\kap_s = \as/4$ relation modifies the \emph{response} of $\psi$ to the
gluon field, i.e., whether it introduces a direct $\psi$--gluon coupling
beyond the gravitational channel.

If $\kap_s/\kap_{\mathrm{EM}} = \as/\alphafine \approx 41$, this provides an
additional enhancement factor of $\sim 40$ for the QCD channel over the EM
channel.

\textbf{Potential contribution: $\times 40$ (if $\kap = \alpha_i/4$ extends
to QCD).}


%=============================================================================
%=============================================================================
\section{Gap-Closing Budget}\label{sec:budget}
%=============================================================================
%=============================================================================

\begin{center}
\renewcommand{\arraystretch}{1.4}
\begin{tabular}{lccc}
\toprule
\textbf{Contribution} & \textbf{Factor} & \textbf{Confidence} &
\textbf{Cumulative} \\
\midrule
Perturbative QCD chain & $\times 20$ & High & $\times 20$ \\
Non-perturbative (instanton + condensate) & $\times 2$--$5$ & Moderate &
$\times 40$--$100$ \\
QCD condensate coherence & $\times 10$--$100$ & Low & $\times 400$--$10^4$ \\
$\kap_s = \as/4$ extension & $\times 40$ & Moderate & $\times 1.6 \times
10^4$--$4 \times 10^5$ \\
Vacuum screening (nuclear) & $\times 10$--$100$ & Low & $\times 1.6 \times
10^5$--$4 \times 10^7$ \\
\midrule
\textbf{Required total} & $\times 2.6 \times 10^5$ & --- & --- \\
\bottomrule
\end{tabular}
\end{center}

The required enhancement of $\sim 2.6 \times 10^5$ over the EM channel is
within the range of the combined mechanisms, but only with optimistic values
for the uncertain factors (condensate coherence and vacuum screening).

\begin{result}[Gap-closing assessment]
The perturbative QCD vertex, combined with established non-perturbative
corrections, provides $\sim 40$--$100 \times$ enhancement over the
EM-only prediction. This leaves a gap of $\sim 2600$--$6500$ (or
$10^{3.4}$--$10^{3.8}$).

With speculative but physically motivated mechanisms (condensate coherence,
vacuum screening, $\kap_s = \as/4$), the gap can be closed, but requires
two or more of these mechanisms operating at their optimistic values.

The UV completion hypothesis ($\ks \sim 1$) closes the gap immediately but
is not derived from established DFD principles.
\end{result}


%=============================================================================
%=============================================================================
\section{Implications for Propulsion}\label{sec:propulsion}
%=============================================================================
%=============================================================================

The propulsion question reduces to: can an engineered electromagnetic system
create a large enough $\dpsi$ to produce a useful force?

The EM-to-$\psi$ coupling goes through:
\begin{equation}
\nabla^2\psi = \frac{2G}{c^4} u_{\mathrm{EM}} + \ldots \, ,
\end{equation}
giving the ``$G/c^4$ wall'': $2G/c^4 \approx 1.6 \times 10^{-44}$~m/J.

The QCD vertex does not directly help propulsion because:
\begin{enumerate}
\item The QCD vertex describes how $\psi$ \emph{affects} QCD, not how QCD
\emph{sources} $\psi$.
\item QCD confinement energy is gravitationally trapped inside nucleons---it
already produces the normal gravitational mass.
\item To source $\psi$ beyond the gravitational channel, one needs an
electromagnetic (or other controllable) source.
\end{enumerate}

However, the QCD vertex is indirectly relevant because it determines the
\emph{sensitivity} of matter to $\psi$ perturbations. If one can create a
$\dpsi$ of order $10^{-6}$ (rather than the gravitational $10^{-10}$), the
QCD vertex tells us the resulting force:
\begin{equation}
F = V_N \cdot (\nabla\dpsi) \sim 28~\text{MeV}/c^2 \cdot \frac{\dpsi}{L}
\sim 5 \times 10^{-12}~\text{N per nucleon} \cdot
\frac{\dpsi}{10^{-6}} \cdot \frac{1~\text{m}}{L} \, .
\end{equation}

For a $1$~kg mass ($\sim 6 \times 10^{26}$ nucleons):
\begin{equation}
F \sim 3 \times 10^{15}~\text{N} \cdot \frac{\dpsi}{10^{-6}} \cdot
\frac{1~\text{m}}{L} \, .
\end{equation}

This is enormous---much larger than gravity ($\sim 10$~N)---but creating
$\dpsi \sim 10^{-6}$ electromagnetically requires overcoming the $G/c^4$
wall, which brings us back to the original problem.

\begin{tcolorbox}[colback=yellow!5!white, colframe=yellow!50!black,
  title=Propulsion Assessment]
The QCD vertex calculation shows that \emph{if} a sufficiently large
$\dpsi$ can be generated, the response of matter through the QCD channel
is $\sim 3500$ times stronger than through the EM channel alone. But
generating $\dpsi$ still requires overcoming the $G/c^4$ wall unless
a non-gravitational sourcing mechanism for $\psi$ is identified.

The $\alpha$-tension data, if confirmed, would demonstrate that either:
(a)~the $G/c^4$ wall has a loophole (e.g., through nuclear resonance), or
(b)~$\psi$ has sources beyond gravity at nuclear scales.
Either conclusion would have revolutionary implications for propulsion.
\end{tcolorbox}


%=============================================================================
%=============================================================================
\section{Honest Assessment}\label{sec:honest}
%=============================================================================
%=============================================================================

%-----------------------------------------------------------------------------
\subsection{What is firmly established}
%-----------------------------------------------------------------------------

\begin{enumerate}
\item \textbf{The gauge-coupling hierarchy.} The DFD framework gives
$\ks/\ka = (\as/\alphafine)^2 \approx 260$ at $M_Z$ and
$\ks/\ka \sim 1600$ at $\mu = 2$~GeV. The strong force is far more
sensitive to $\psi$ than electromagnetism.

\item \textbf{The coupling chain.} The chain
$\psi \to \as \to \LQCD \to \mN$ is well-defined and calculable in
perturbation theory. The result is $V_{\mathrm{QCD}}^{(\mathrm{pert})}
\approx 28~\text{MeV} \cdot \dpsi$.

\item \textbf{The species-dependent mechanism.} The $\fEM$ proportionality
arises naturally from the interplay between QCD-shifted nucleon mass and
fixed Coulomb energy.

\item \textbf{The perturbative gap.} Even with the QCD enhancement, the
perturbative prediction falls short of the $\alpha$-tension data by a factor
of $\sim 260$. This gap is robust.
\end{enumerate}

%-----------------------------------------------------------------------------
\subsection{What is uncertain}
%-----------------------------------------------------------------------------

\begin{enumerate}
\item \textbf{Non-perturbative QCD.} The instanton and gluon condensate
contributions are estimated at $\sim 2$--$5\times$ enhancement, but with
large uncertainties. Lattice QCD could provide a definitive answer if
simulations with modified $\as$ were performed.

\item \textbf{Condensate coherence.} The claim that the QCD condensate
enhances the coupling by a factor of $\sim 10$--$100$ through coherence
effects is speculative. Unlike Cooper pairs, the QCD condensate does not
have a macroscopic wavefunction with a well-defined phase.

\item \textbf{The $\kap = \alpha_i/4$ extension.} The vacuum loading
derivation has been performed only for U(1). Its extension to SU(3) is a
conjecture that requires a separate derivation from the non-Abelian
gauge-emergence framework.

\item \textbf{Vacuum screening at nuclear scales.} The nuclear resonance
mechanism is physically motivated but the enhancement factors depend on
the GDR--$\psi$ coupling, which is not well constrained.
\end{enumerate}

%-----------------------------------------------------------------------------
\subsection{What would close the gap}
%-----------------------------------------------------------------------------

Three paths forward:

\paragraph{Path A: Establish the $\kap_s = \as/4$ relation.}
If the vacuum loading mechanism extends to SU(3), this provides $\times 40$
enhancement. Combined with the perturbative chain ($\times 20$) and modest
non-perturbative corrections ($\times 3$), this gives $\times 2400$---close
to the required $\times 2600$. This is the most calculable path and should
be the next priority.

\paragraph{Path B: Lattice QCD calibration.}
A direct lattice QCD calculation of $\partial \mN/\partial \as$ at the
non-perturbative level would resolve the instanton/condensate uncertainties.
If the lattice result gives $\partial\ln\mN/\partial\ln\as \sim 0.5$--$1$
(rather than the perturbative $\sim 0.03$), the gap would be largely closed
within the established framework.

\paragraph{Path C: UV completion.}
If the DFD $\mufunction(\psi)$ has a UV branch with confining dynamics at
the fm scale, the coupling becomes $\order{1}$ and the gap closes trivially.
This is the most elegant solution but the hardest to derive.

%-----------------------------------------------------------------------------
\subsection{Summary numbers}\label{sec:summary_numbers}
%-----------------------------------------------------------------------------

\begin{center}
\renewcommand{\arraystretch}{1.3}
\begin{tabular}{lcc}
\toprule
\textbf{Quantity} & \textbf{Value} & \textbf{Units} \\
\midrule
$\ka$ (U(1) gauge-coupling coeff.) & $8.5 \times 10^{-6}$ & --- \\
$\ks$ at $\mu = 2$~GeV & $1.4 \times 10^{-2}$ & --- \\
$\ks$ at $\mu = M_Z$ & $2.2 \times 10^{-3}$ & --- \\
$\mathcal{A}_\Lambda$ at $\mu = 2$~GeV & 2.3 & --- \\
$f_\Lambda$ & 0.91 & --- \\
$V_{\mathrm{QCD}}^{(\mathrm{pert})}$ & 28 & MeV $\cdot \dpsi$ \\
$V_{\mathrm{QCD}}^{(\mathrm{best})}$ & 60 (range: 28--140) & MeV $\cdot \dpsi$ \\
$C_{\mathrm{th}}$ & $2.7 \times 10^{-2}$ & --- \\
$C_{\mathrm{obs}}$ (from $\alpha$-tension) & $1.765 \times 10^{-6}$ & --- \\
Gap: $C_{\mathrm{obs}} / (C_{\mathrm{th}} \cdot \dpsi_\oplus)$ & $\sim 260$ & --- \\
Required enhancement beyond pert. & $50$--$260$ & --- \\
\bottomrule
\end{tabular}
\end{center}


%=============================================================================
%=============================================================================
\section{Conclusions}\label{sec:conclusions}
%=============================================================================
%=============================================================================

We have computed the $\psi$--QCD vertex from first principles within the DFD
gauge-emergence framework. The key results are:

\begin{enumerate}
\item The SU(3) gauge-coupling coefficient $\ks = \as^2/(2\pi)$ makes the
strong force $\sim 260$--$1600$ times more sensitive to $\psi$ than
electromagnetism, depending on the renormalization scale.

\item The perturbative coupling chain $\psi \to \as \to \LQCD \to \mN$ gives
a $\psi$--nucleon vertex of $V_{\mathrm{QCD}} \approx 28$~MeV per unit
$\dpsi$, or about $3\%$ of the nucleon mass.

\item This perturbative result leaves a gap of $\sim 10^{2.4}$ relative to
the $\alpha$-tension data. The gap is large but \emph{dramatically smaller}
than the $\sim 10^{22}$ gap from the gravitational estimate or the
$\sim 10^{5}$ gap from the EM-only estimate.

\item Non-perturbative QCD effects (instantons, gluon condensate, string
tension) enhance the vertex by a factor of $\sim 2$--$5$, reducing the gap
to $\sim 50$--$130$.

\item Three candidate mechanisms for closing the remaining gap are identified.
The most tractable is the extension of the vacuum loading relation
$\kap = \alphafine/4$ to the strong sector ($\kap_s = \as/4$), which would
provide a factor of $\sim 40$ and bring the prediction within striking
distance of the observation.

\item The $\psi$--QCD vertex does not directly help propulsion (the $G/c^4$
wall remains for sourcing $\psi$), but it shows that the response of matter
to $\psi$ perturbations is $\sim 3500$ times stronger through the QCD channel
than the EM channel.
\end{enumerate}

The most important next calculations for the DFD programme are:
\begin{enumerate}
\item Derive (or refute) $\kap_s = \as/4$ from the non-Abelian gauge-emergence
framework.
\item Obtain the lattice QCD value of $\partial\ln\mN/\partial\ln\as$ to
anchor the non-perturbative enhancement factor.
\item Investigate the UV completion of $\mufunction(\psi)$ at the fm scale.
\end{enumerate}


%=============================================================================
% REFERENCES
%=============================================================================
\begin{thebibliography}{99}

\bibitem{DFD_Unified}
G.~T.~Alcock,
``Density Field Dynamics: A Complete Unified Theory,''
DFD Research Programme (2025);
see Section~8B and Appendix~F.

\bibitem{Ji1995}
X.~Ji,
``A QCD analysis of the mass structure of the nucleon,''
Phys.\ Rev.\ Lett.\ \textbf{74}, 1071 (1995).

\bibitem{Yang2018}
Y.-B.~Yang \emph{et al.} ($\chi$QCD Collaboration),
``Proton mass decomposition from the QCD energy momentum tensor,''
Phys.\ Rev.\ Lett.\ \textbf{121}, 212001 (2018).

\bibitem{SVZ1979}
M.~A.~Shifman, A.~I.~Vainshtein, and V.~I.~Zakharov,
``QCD and resonance physics: theoretical foundations,''
Nucl.\ Phys.\ B \textbf{147}, 385 (1979).

\bibitem{Damour2010}
T.~Damour and J.~F.~Donoghue,
``Equivalence principle violations and couplings of a light dilaton,''
Phys.\ Rev.\ D \textbf{82}, 084033 (2010);
arXiv:1007.2792 [gr-qc].

\bibitem{QCD_Psi}
G.~T.~Alcock,
``Does the DFD $\psi$-Field Couple to QCD Confinement Energy?
The Magnitude Gap Problem and the Strong-Sector Amplification Channel,''
DFD Research Programme (2026).

\bibitem{VacScreen}
G.~T.~Alcock,
``Vacuum Screening of the Gauge Coupling in Density Field Dynamics:
Species-Dependent $\alpha$ Shifts Without Equivalence-Principle Violation,''
DFD Research Programme (2026).

\bibitem{NucRes}
G.~T.~Alcock,
``Nuclear Resonance Vacuum Screening in Density Field Dynamics:
Giant Dipole Resonance Coupling to the $\psi$ Field,''
DFD Research Programme (2026).

\bibitem{CoreField}
G.~T.~Alcock,
``Core Derivations of the DFD $\psi$-Field Equation
with Electromagnetic Sources, Quantum Coherence Enhancement,
Self-Consistent Back-Reaction, and Directed Force Generation,''
DFD Research Programme (2025).

\bibitem{VacLoading}
G.~T.~Alcock,
``Vacuum Loading of the $\psi$ Field: Derivation of $\kappa = \alpha/4$
from Density Field Dynamics,''
DFD Research Programme (2026).

\bibitem{Shuryak1982}
E.~V.~Shuryak,
``The role of instantons in quantum chromodynamics,''
Nucl.\ Phys.\ B \textbf{203}, 93 (1982).

\bibitem{Diakonov1986}
D.~Diakonov and V.~Petrov,
``Instanton-based vacuum from the Feynman variational principle,''
Nucl.\ Phys.\ B \textbf{272}, 457 (1986).

\bibitem{MICROSCOPE2022}
P.~Touboul \emph{et al.} (MICROSCOPE Collaboration),
``MICROSCOPE mission: final results of the test of the equivalence
principle,''
Phys.\ Rev.\ Lett.\ \textbf{129}, 121102 (2022).

\bibitem{FLAG2021}
Y.~Aoki \emph{et al.} (FLAG Working Group),
``FLAG Review 2021,''
Eur.\ Phys.\ J.\ C \textbf{82}, 869 (2022).

\bibitem{Hatta2018}
Y.~Hatta, A.~Rajan, and D.-L.~Yang,
``Quark and gluon contributions to the QCD trace anomaly,''
Phys.\ Rev.\ D \textbf{98}, 074003 (2018).

\bibitem{Lorce2021}
C.~Lorc\'{e}, A.~Metz, B.~Pasquini, and S.~Rodini,
``Energy-momentum tensor in QCD: nucleon mass decomposition and mechanical
equilibrium,''
JHEP \textbf{11}, 121 (2021).

\end{thebibliography}

\end{document}
